%% =========================================================
%% Birkhoff's theorem --- doubly-stochastic half
%% =========================================================

\section{Birkhoff's theorem}

Birkhoff's theorem characterizes the convex geometry of doubly
stochastic matrices.

\begin{theorem}[Birkhoff (doubly-stochastic part)]
    \label{thm:birkhoff_doubly_stochastic}
    \lean{TNLean.wolf_birkhoff_doublyStochastic_convexHull}
    \lean{TNLean.wolf_birkhoff_extremePoints_doublyStochastic}
    \leanok
    The set of doubly stochastic matrices in $M_d(\R)$ is the convex
    hull of the $d\times d$ permutation matrices, and its extreme
    points are exactly the permutation matrices.
    See \cite[Chapter~8, Theorem~8.6]{Wolf2012Quantum}.
    This statement concerns only doubly stochastic matrices; the
    corresponding assertion of Theorem~8.6 for doubly substochastic
    matrices is not included here.
\end{theorem}
% The doubly-substochastic half of Birkhoff's theorem (extreme
% points of the substochastic matrices) is not formalized; #3959.

\begin{proof}\leanok
    Let $\mathcal{DS}_d \subset M_d(\mathbb{R})$ be the set of $d \times d$
    doubly stochastic matrices and let $\mathcal{P}_d$ be the set of
    $d \times d$ permutation matrices.  The Birkhoff--von Neumann
    decomposition gives
    \begin{align}
      \mathcal{DS}_d &= \operatorname{conv}(\mathcal{P}_d), \\
      \operatorname{ext}(\mathcal{DS}_d) &= \mathcal{P}_d .
    \end{align}
    Indeed, every doubly stochastic matrix is a convex combination of
    permutation matrices, while convexity of $\mathcal{DS}_d$ gives the
    converse inclusion in the first identity.  Suppose that a permutation
    matrix $P_\sigma$ is a convex combination
    $P_\sigma=tA+(1-t)B$, where $0<t<1$ and $A,B$ are doubly stochastic.
    Nonnegativity and the row sums give
    \begin{align}
      (P_\sigma)_{ij}=0
      &\Longrightarrow A_{ij}=B_{ij}=0, \\
      1=\sum_j A_{ij} &= A_{i,\sigma(i)}, \\
      1=\sum_j B_{ij} &= B_{i,\sigma(i)}.
    \end{align}
    Thus $A=B=P_\sigma$, so every permutation matrix is extreme.
    Conversely, a non-permutation matrix has a Birkhoff decomposition
    involving at least two distinct permutation matrices and is therefore
    not extreme, which proves the second identity.
\end{proof}

\section{Entropy increasing maps}

\begin{theorem}[Concave traces under unital positive maps]
    \label{thm:concave_trace_positive_map}
    \lean{IsPositiveMap.trace_concave_le}
    \leanok
    Let $T:M_d(\mathbb C)\to M_d(\mathbb C)$ be positive, unital, and
    trace preserving, and let $h:\mathbb R\to\mathbb R$ be concave.
    Every positive semidefinite matrix $A$ satisfies
    $\operatorname{tr}h(A)\le\operatorname{tr}h(T(A))$.
    This is the forward implication of \cite[Equation~(8.22)]{Wolf2012Quantum},
    without trace normalization or strict concavity.
\end{theorem}
\begin{proof}\leanok
    Choose unitary matrices $U,V$ diagonalizing $A,T(A)$, with eigenvalues
    $\lambda_j,\mu_i$, and set
    $S(X)=V^*T(UXU^*)V$ and $M_{ij}=(S(E_{jj}))_{ii}$.
    Positivity gives $M_{ij}\ge0$, unitality gives $\sum_jM_{ij}=1$,
    and trace preservation gives $\sum_iM_{ij}=1$.
    Moreover, $\mu_i=\sum_jM_{ij}\lambda_j$.
    Scalar Jensen and the column sums yield
    $\sum_jh(\lambda_j)\le\sum_ih(\mu_i)$.
    The spectral formula for the functional calculus identifies these sums
    with the two traces.
\end{proof}

\begin{theorem}[Uniform bound for concave traces]
    \label{thm:concave_trace_uniform}
    \lean{Matrix.IsHermitian.trace_concave_le_uniform}
    \leanok
    For $d>0$, a Hermitian matrix $A\in M_d(\mathbb C)$ with trace one,
    and a concave function $h:\mathbb R\to\mathbb R$,
    $\operatorname{tr}h(A)\le d\,h(1/d)$.
\end{theorem}
\begin{proof}\leanok
    The eigenvalues sum to one. Apply scalar Jensen with the positive
    uniform weights $1/d$, and multiply by $d$.
\end{proof}

\begin{theorem}[Strict concavity and the uniform maximizer]
    \label{thm:strict_concave_trace_uniform}
    \lean{Matrix.IsHermitian.eq_smul_one_of_trace_strictConcave_ge}
    \leanok
    For $d>0$, let $A\in M_d(\mathbb C)$ be Hermitian with trace one,
    and let $h:\mathbb R\to\mathbb R$ be strictly concave.
    If $d\,h(1/d)\le\operatorname{tr}h(A)$, then $A=I/d$.
\end{theorem}
\begin{proof}\leanok
    The assumed inequality reverses uniform Jensen. Its strict equality
    criterion makes all eigenvalues equal. Their sum is one, so each is
    $1/d$. The spectral decomposition then gives $A=I/d$.
\end{proof}

\begin{theorem}[Entropy increasing maps]
    \label{thm:entropy_increasing_maps}
    \lean{IsPositiveMap.unital_iff_trace_strictConcave_increasing}
    \leanok
    Let $T:M_d(\mathbb C)\to M_d(\mathbb C)$ be positive and trace
    preserving, and let $h:\mathbb R\to\mathbb R$ be strictly concave.
    Then $T(I)=I$ if and only if
    $\operatorname{tr}h(\rho)\le\operatorname{tr}h(T(\rho))$
    for every density matrix $\rho$.
    This is \cite[Equation~(8.22), Corollary ``Entropy increasing maps'']{Wolf2012Quantum}.
\end{theorem}
\begin{proof}\leanok
    \uses{thm:concave_trace_positive_map,thm:strict_concave_trace_uniform}
    The forward implication follows from
    Theorem~\ref{thm:concave_trace_positive_map}.
    For the converse, suppose $d>0$ and take $\rho=I/d$.
    Its image is positive semidefinite with trace one, and the assumed
    entropy increase gives
    $d\,h(1/d)\le\operatorname{tr}h(T(I/d))$.
    Theorem~\ref{thm:strict_concave_trace_uniform} gives $T(I/d)=I/d$;
    linearity implies $T(I)=I$.
    When $d=0$, all matrices are equal and unitality is automatic.
\end{proof}
\endinput
